% Generated by task tex-docs from dossiers/docs-debt.md, sections E (who owns what) and F (the codes).
% Do not edit by hand.
\section{Proposal: who owns what}\label{sec:docs-proposal-owners}

\subsection{The rule}\label{sec:docs-own-rule}

A fact has one home. Every other place links to the home and restates nothing: no summary, no ``for convenience'' copy, no table repeated with fewer columns. A reader who wants the fact follows the link; a writer who changes the fact changes one place.

\subsection{The owners}\label{sec:docs-owners}

{\footnotesize\setlength{\tabcolsep}{3pt}
\begin{longtable}{@{}>{\raggedright\arraybackslash}p{0.291\linewidth}>{\raggedright\arraybackslash}p{0.233\linewidth}>{\raggedright\arraybackslash}p{0.446\linewidth}@{}}
\caption{One home per kind of fact.}\label{tab:docs-owners}\\
\toprule
\textbf{Kind of fact} & \textbf{Owner} & \textbf{What the other places do} \\
\midrule
\endfirsthead
\toprule
\textbf{Kind of fact} & \textbf{Owner} & \textbf{What the other places do} \\
\midrule
\endhead
\bottomrule
\endfoot
Targets, scope, conventions, the spine, acceptance tests & the blueprint & the status and the tracker link to the section \\\addlinespace[2pt]
The specification of each unit of work: what it produces, with signatures; what it needs; its sources; its tests; the upstream work to read first & the blueprint, one section per unit & the hole comment is reduced to a pointer; the checklist line carries a name and a link, no signature \\\addlinespace[2pt]
Decisions, taken and open & the blueprint, a section ``Decisions'' & the status and the tracker hold none \\\addlinespace[2pt]
The upstream ledger: what is admitted or absent at the pin, what replaces it here, the upstream number, when the local copy goes & the blueprint, inside ``Dependencies and exact contracts'' & the status's and the tracker's copies are deleted \\\addlinespace[2pt]
Coverage: what is on \texttt{main} & the repository; the status prints it, and a script checks the print & the tracker's checklist is ticked from it; the blueprint says nothing \\\addlinespace[2pt]
What can start now & the status, derived from the units' ``needs'' and the coverage & the tracker links \\\addlinespace[2pt]
Upstream watch: external work in flight, with its adoption condition & the status & the blueprint's column ``Existing work'' is deleted; its content becomes the ``upstream work to read first'' of the unit, without state \\\addlinespace[2pt]
Who is doing what & the tracker's checklist line, linking the claimant's intention issue or pull request & the status holds no claim \\\addlinespace[2pt]
Open pull requests & GitHub & no table anywhere \\\addlinespace[2pt]
Discrepancies in the leanVM sources at the pin (specification, Rust, Python) & \texttt{docs/\allowbreak{}leanvm-\allowbreak{}target.\allowbreak{}md}, a new section ``Known discrepancies at the pin'', for both roadmaps & the blueprint cites a discrepancy by name where it binds a definition; both status files lose the section \\\addlinespace[2pt]
Limits of the pinned libraries & \texttt{docs/\allowbreak{}dependencies.\allowbreak{}md}, under each library & the same \\\addlinespace[2pt]
Lean pitfalls & \texttt{tests/\allowbreak{}README.\allowbreak{}md}, which already holds five of them, linked from \texttt{docs/\allowbreak{}development.\allowbreak{}md} & the same \\\addlinespace[2pt]
The record of a review & the handoff under \texttt{docs/\allowbreak{}reviews/\allowbreak{}} & see \cref{sec:docs-handoffs} \\\addlinespace[2pt]
What was read, and when & an archive (\cref{sec:docs-survey-record}) & - \\\addlinespace[2pt]
How to work on the proof system & the blueprint's section ``How work is tracked''; the general rules stay in \texttt{CONTRIBUTING.\allowbreak{}md} and \texttt{AGENTS.\allowbreak{}md} & the tracker links \\
\end{longtable}}

No new kind of document is created. Three existing documents each gain one section (\texttt{leanvm-\allowbreak{}target.\allowbreak{}md}, \texttt{dependencies.\allowbreak{}md}, \texttt{tests/\allowbreak{}README.\allowbreak{}md}), which is where their subject already lives: \texttt{dependencies.\allowbreak{}md} already explains two limits of Clean and cites a finding by its code (\texttt{:\allowbreak{}52-\allowbreak{}55}, \texttt{:\allowbreak{}81-\allowbreak{}88}), and \texttt{tests/\allowbreak{}README.\allowbreak{}md} already explains five pitfalls.

Why \texttt{leanvm-\allowbreak{}target.\allowbreak{}md} for the discrepancies. They are facts about the target at the pin, not about either roadmap: they hold until the pin moves, they are shared (the missing section 8.4 is one finding in both status files), and both status files already say they are ``numbered for citation from ... \texttt{docs/\allowbreak{}leanvm-\allowbreak{}target.\allowbreak{}md}''. One register also ends the collision between the two files' letters. The alternative, an appendix of the blueprint, keeps two registers, one per roadmap, and is not recommended.

\subsection{What the blueprint contains, and what leaves it}\label{sec:docs-blueprint-contents}

Table of contents proposed:

\begin{enumerate}
\item What is proved: the two master theorems, the compiled verifier's theorem, the two base theorems (T4) (today's introduction).
\item For zkVM engineers.
\item Scope.
\item Dependencies and exact contracts, ending with the upstream ledger (one table).
\item Pinned conventions.
\item The spine.
\item The units of work: one table with names (\cref{sec:docs-text-holes}), then one section per unit in build order. Each section merges today's per-layer sketch and the hole comment's section, and has the same five parts: produces, needs, sources, tests, upstream work to read first.
\item Acceptance tests and nearby false statements.
\item Interfaces supplied to later work.
\item Boundaries with the other roadmaps.
\item Decisions: a table of those taken (number, name, the choice in one line, the convention or test that records it) and a list of those open.
\item How work is tracked (\cref{sec:docs-text-tracking}).
\item Index of names and numbers (\cref{sec:docs-codes-index-draft}).
\item References.
\end{enumerate}

What leaves the blueprint: every sentence of status or history listed in \cref{sec:docs-sketches}, item 11; the columns ``Existing work'' and ``Issue'' of the hole table; the section ``Ordering and parallelism'', whose content is the column ``needs''; the per-layer sketches that contradict the spine, which are rewritten on the spine's types when merged into section 7.

What enters it: the tests and upstream notes that exist only in the hole comment (end of \cref{sec:docs-inventory}); the statement of decisions 1 to 10, of which only the outcomes are written today; what the spine asks of the Flock roadmap, which is a comment on issue 3.

\subsection{What the status contains, and what can be generated}\label{sec:docs-status-contents}

Table of contents proposed:

\begin{enumerate}
\item Head: two sentences. What the file is; that the coverage table is checked by a script.
\item Coverage: one row per unit, generated.
\item What can start now: generated, or five lines by hand.
\item Upstream watch: the item, the unit, what it would replace, the adoption condition. No state and no date: the state is one click away and is what goes stale.
\item Pointers: to the blueprint's decisions and ledger, to the three registers, to the reviews.
\end{enumerate}

\textbf{Coverage from the repository.} The repository forbids \texttt{sorry}, and the kernel axiom audit rejects \texttt{sorryAx}, so a declaration that exists is proved. A unit is therefore built exactly when the declarations that mark it exist. The probe \texttt{coverage\_\allowbreak{}from\_\allowbreak{}blueprint.\allowbreak{}py} does this today from the blueprint's hole table, with no Lean: it reports the tables-and-stacking unit at 15 of 15 names, the composition at 2 of 2, the public-input phase at 1 of 1, and every unbuilt unit at 0, except where the table lists names that are not the unit's own (\cref{sec:docs-repository}, item 6). What it needs from the blueprint is a column ``built when declared'' holding, for each unit, names that this unit and no other declares. The public-input phase sets the pattern: a phase \texttt{x} is marked by \texttt{xPhase} and \texttt{xComplete} for its first half and \texttt{xSecurity} for its second.

Proposed, for the owner to decide: \texttt{scripts/\allowbreak{}check-\allowbreak{}protocol-\allowbreak{}status.\allowbreak{}py}, run by \texttt{scripts/\allowbreak{}validate.\allowbreak{}sh} and by the workflow ``Documentation integrity''. It reads the units table of the blueprint, looks each marking name up under \texttt{LeanerVM/\allowbreak{}}, and compares the result with the table between two markers in the status; on a difference it fails and prints the table to paste, and with \texttt{-\allowbreak{}-\allowbreak{}write} it rewrites it. A planted difference goes into \texttt{scripts/\allowbreak{}test-\allowbreak{}policy-\allowbreak{}checks.\allowbreak{}py}, as for the other policy scripts. The same script can check two more things this review found by hand: that every name of the interface list that belongs to a built unit is declared, and the import rule of the proof system (no \texttt{import Leaner\allowbreak{}VM.\allowbreak{}Arithmetization} under \texttt{Leaner\allowbreak{}VM/\allowbreak{}Protocol/\allowbreak{}} outside an allow-list), which the Layer 1 review left open.

The head then needs no commit hash. A file cannot name the commit that contains it; today's head names the parent and is wrong by construction after every merge (\cref{sec:docs-status-stale}, items 1 to 5).

\textbf{Upstream state.} \texttt{scripts/\allowbreak{}check-\allowbreak{}upstreams.\allowbreak{}sh} already needs \texttt{gh} and the network and runs weekly in the workflow ``Upstream drift''. It can print the state of each watched item next to the pins' drift. The state then lives in the workflow's report and is never committed.

\subsection{The survey record}\label{sec:docs-survey-record}

The 190 lines of ``Survey record'' (\texttt{protocol-\allowbreak{}status.\allowbreak{}md:\allowbreak{}434-\allowbreak{}622}) are a log of what was read. They are history, useful ``so the searches are not repeated'', and they are what makes the status long. Two options. Recommended: move them once, unchanged, to \texttt{docs/\allowbreak{}reviews/\allowbreak{}protocol-\allowbreak{}survey-\allowbreak{}record.\allowbreak{}md}, headed as an archive; the directory is already the repository's archive of reviews. Alternative: keep them as the last section of the status, headed ``Record (history, not a source of truth)'', which keeps a hand-maintained status long.

\subsection{The review handoffs}\label{sec:docs-handoffs}

A handoff is a record of what was found at one commit. It is never a specification.

\begin{itemize}
\item A finding that is accepted becomes blueprint text (a convention, a signature, an acceptance test, a decision) in the pull request that meets it. The handoff's disposition table cites the blueprint's section, by name.
\item The handoff gets one line at its head: the commit it describes, and ``an archive; names and paths are the reviewed commit's; for the current state read the blueprint''.
\item The status does not summarise reviews. \texttt{docs/\allowbreak{}README.\allowbreak{}md} keeps its index of handoffs, one line each, without dispositions.
\item Texts a handoff drafts for GitHub (\texttt{protocol-\allowbreak{}spine.\allowbreak{}md:\allowbreak{}662-\allowbreak{}769}) are applied by editing the blueprint, since the specifications they patch move there.
\end{itemize}

\subsection{The tracker and the hole comment}\label{sec:docs-tracker-proposal}

The body of issue 12: the pointers, the checklist, and two lines on how to claim (\cref{sec:docs-text-tracker}). No ledger, no watch list, no table of pull requests, no decisions, no list of names, no ``frontier''.

The hole comment: its thirteen sections move to the blueprint, corrected on the way (\cref{sec:docs-hole-comment}). The comment is then edited by its author to one line, ``The specifications of the units of work are in the blueprint, section 'The units of work''', so that links to it still land somewhere.

The three comments by the other contributor are a record of claims made in September and need nothing. Their content that is still live (four claims) is already in the bodies of intention issues 28, 31, 33, 37.

\textbf{Claiming without comments.} A maintainer, or an agent acting for one, edits the checklist line in the body. A contributor without write access cannot edit the body; they open an intention issue whose body names the unit and the exact declarations taken, which they can edit themselves, and the maintainer links it from the checklist line. Nobody claims by comment. This is what the four live claims already do.

\section{Proposal: the codes}\label{sec:docs-proposal-codes}

\subsection{What is replaced by names, and what stays}\label{sec:docs-codes-fate}

{\footnotesize\setlength{\tabcolsep}{3pt}
\begin{longtable}{@{}>{\raggedright\arraybackslash}p{0.236\linewidth}>{\raggedright\arraybackslash}p{0.747\linewidth}@{}}
\caption{What is replaced by names, and what stays.}\label{tab:docs-codes-fate}\\
\toprule
\textbf{Family} & \textbf{Proposal} \\
\midrule
\endfirsthead
\toprule
\textbf{Family} & \textbf{Proposal} \\
\midrule
\endhead
\bottomrule
\endfoot
Units of work & names. The letters \texttt{S}, \texttt{G}, \texttt{I}, \texttt{P}, \texttt{C}, \texttt{K}, \texttt{L} go \\\addlinespace[2pt]
Ledger entries & names \\\addlinespace[2pt]
Findings of reviews & names, in the handoff's own disposition table; never cited by letter outside the handoff \\\addlinespace[2pt]
Findings against the leanVM sources, the libraries, Lean & names, each a heading in its register, so that it has an anchor; the former code is kept once, in the register, to read old pull requests \\\addlinespace[2pt]
Layers 0 to 13 & stay, as the numbers of the blueprint's sections, with the name: ``the bus phase (Layer 6)'' \\\addlinespace[2pt]
Acceptance tests 1 to 28 & stay, with the name the blueprint already gives each in bold \\\addlinespace[2pt]
Decisions 1 to 15 & stay, each given a name \\\addlinespace[2pt]
Target theorems T1 to T8 & stay; owned by \texttt{architecture.\allowbreak{}md} \\\addlinespace[2pt]
Category A, Category B & replaced by the two words they stand for: ``written from the specification'', ``transcribed from its source'' \\
\end{longtable}}

The rule for what stays: \textbf{the name stands beside the number at every use.} ``Acceptance test 24'' alone is not allowed; ``the extractor computes (acceptance test 24)'' is. A reference to the other roadmap's layer, test or decision always says ``leanISA''.

If the owner prefers to keep numbers for the findings against the sources, the scheme without collisions is a single counter in \texttt{docs/\allowbreak{}leanvm-\allowbreak{}target.\allowbreak{}md}, without letters, across both roadmaps (``source finding 31''), with the former codes in one column. Letters are what collide.

\subsection{The mapping from every code to its name}\label{sec:docs-codes-mapping}

\textbf{Units of work.} The short form is for branch names, pull-request titles and the checklist.

{\footnotesize\setlength{\tabcolsep}{3pt}
\begin{longtable}{@{}>{\raggedright\arraybackslash}p{0.116\linewidth}>{\raggedright\arraybackslash}p{0.582\linewidth}>{\raggedright\arraybackslash}p{0.271\linewidth}@{}}
\caption{The units of work: former code, name, short form.}\label{tab:docs-map-units}\\
\toprule
\textbf{Former code} & \textbf{Name} & \textbf{Short form} \\
\midrule
\endfirsthead
\toprule
\textbf{Former code} & \textbf{Name} & \textbf{Short form} \\
\midrule
\endhead
\bottomrule
\endfoot
\texttt{0} & field instances (Layer 0) & \texttt{field-\allowbreak{}instances} \\\addlinespace[2pt]
\texttt{L1} & tables and stacking (Layer 1) & \texttt{tables-\allowbreak{}stacking} \\\addlinespace[2pt]
\texttt{S} & the spine & \texttt{spine} \\\addlinespace[2pt]
\texttt{C1} & knowledge-soundness composition (the port of ArkLib pull request 615) & \texttt{knowledge-\allowbreak{}append} \\\addlinespace[2pt]
\texttt{I1} & Clean expressions as polynomials (Layer 2) & \texttt{clean-\allowbreak{}polynomials} \\\addlinespace[2pt]
\texttt{I2} & the adaptor (Layer 3) & \texttt{adaptor} \\\addlinespace[2pt]
\texttt{G1} & sumcheck: definition and completeness (Layer 4) & \texttt{sumcheck-\allowbreak{}def} \\\addlinespace[2pt]
\texttt{G2} & sumcheck: knowledge soundness (Layer 4) & \texttt{sumcheck-\allowbreak{}security} \\\addlinespace[2pt]
\texttt{G3} & batching by powers (Layer 4) & \texttt{batching} \\\addlinespace[2pt]
\texttt{G4} & fingerprint and collision bound (Layer 5) & \texttt{fingerprint} \\\addlinespace[2pt]
\texttt{G5} & grand-product GKR: definition and completeness (Layer 5) & \texttt{gkr-\allowbreak{}def} \\\addlinespace[2pt]
\texttt{G6} & grand-product GKR: knowledge soundness (Layer 5) & \texttt{gkr-\allowbreak{}security} \\\addlinespace[2pt]
\texttt{P1} & bus phase: definition and completeness (Layer 6) & \texttt{bus-\allowbreak{}def} \\\addlinespace[2pt]
\texttt{P2} & bus phase: knowledge soundness (Layer 6) & \texttt{bus-\allowbreak{}security} \\\addlinespace[2pt]
\texttt{P3} & table sumcheck phase: definition and completeness (Layer 7) & \texttt{table-\allowbreak{}def} \\\addlinespace[2pt]
\texttt{P4} & table sumcheck phase: knowledge soundness (Layer 7) & \texttt{table-\allowbreak{}security} \\\addlinespace[2pt]
\texttt{P5} & public-input phase (Layer 8) & \texttt{public-\allowbreak{}input} \\\addlinespace[2pt]
\texttt{P6} & Flock phase, the interface to the Flock roadmap (Layer 9) & \texttt{flock-\allowbreak{}interface} \\\addlinespace[2pt]
\texttt{P7} & opening phase: definition and completeness (Layer 10) & \texttt{opening-\allowbreak{}def} \\\addlinespace[2pt]
\texttt{P8} & opening phase: knowledge soundness (Layer 10) & \texttt{opening-\allowbreak{}security} \\\addlinespace[2pt]
\texttt{K1} & WHIR opening (Layer 11) & \texttt{whir} \\\addlinespace[2pt]
\texttt{K2} & Merkle trees, byte hasher, WHIR parameters (Layer 11) & \texttt{merkle-\allowbreak{}parameters} \\\addlinespace[2pt]
\texttt{K3} & transcript, proof object, compiled verifier (Layer 12) & \texttt{compiled-\allowbreak{}verifier} \\\addlinespace[2pt]
\texttt{K4} & the base theorems of T4 (Layer 13) & \texttt{base-\allowbreak{}theorems} \\
\end{longtable}}

\textbf{Ledger entries.}

{\footnotesize\setlength{\tabcolsep}{3pt}
\begin{longtable}{@{}>{\raggedright\arraybackslash}p{0.147\linewidth}>{\raggedright\arraybackslash}p{0.835\linewidth}@{}}
\caption{The upstream ledger: former code and name.}\label{tab:docs-map-ledger}\\
\toprule
\textbf{Former code} & \textbf{Name} \\
\midrule
\endfirsthead
\toprule
\textbf{Former code} & \textbf{Name} \\
\midrule
\endhead
\bottomrule
\endfoot
\texttt{A1} & admitted upstream: sumcheck round knowledge soundness \\\addlinespace[2pt]
\texttt{A2} & admitted upstream: composition of knowledge soundness \\\addlinespace[2pt]
\texttt{A3} & admitted upstream: round-by-round implies plain \\\addlinespace[2pt]
\texttt{A4} & admitted upstream: context lifting (not consumed) \\\addlinespace[2pt]
\texttt{A5} & admitted or absent upstream: Fiat-Shamir and BCS security \\\addlinespace[2pt]
\texttt{A6} & absent upstream: grand product, GKR, batching, stacking \\\addlinespace[2pt]
\texttt{A7} & absent upstream: WHIR, Merkle trees, BLAKE2s \\\addlinespace[2pt]
\texttt{A8} & admitted upstream: correlated agreement up to Johnson \\\addlinespace[2pt]
\texttt{A9} & admitted upstream: ring-switching packing \\\addlinespace[2pt]
\texttt{C1} & absent in Clean: expressions as polynomials \\\addlinespace[2pt]
\texttt{C2} & asked of leanISA: power-of-two heights, bus data per channel (done) \\\addlinespace[2pt]
\texttt{C3} & absent in Clean: balance counted in the naturals \\\addlinespace[2pt]
\texttt{C4} & absent in Clean: fixed columns and sound prover data \\
\end{longtable}}

\textbf{Findings against the leanVM sources} (to \texttt{docs/\allowbreak{}leanvm-\allowbreak{}target.\allowbreak{}md}).

{\footnotesize\setlength{\tabcolsep}{3pt}
\begin{longtable}{@{}>{\raggedright\arraybackslash}p{0.147\linewidth}>{\raggedright\arraybackslash}p{0.835\linewidth}@{}}
\caption{The findings against the leanVM sources: former code and name.}\label{tab:docs-map-sources}\\
\toprule
\textbf{Former code} & \textbf{Name} \\
\midrule
\endfirsthead
\toprule
\textbf{Former code} & \textbf{Name} \\
\midrule
\endhead
\bottomrule
\endfoot
\texttt{S6} & Fiat-Shamir section is \texttt{TODO} \\\addlinespace[2pt]
\texttt{S9} & proof of the product lemma is \texttt{TODO} \\\addlinespace[2pt]
\texttt{S10} & degree of the radix-4 round polynomial not stated \\\addlinespace[2pt]
\texttt{S11} & one bus root not stated \\\addlinespace[2pt]
\texttt{S12} & grinding per level not in Annex B \\\addlinespace[2pt]
\texttt{S13} & retired \\\addlinespace[2pt]
\texttt{S14} & the PDF is not the pinned text \\\addlinespace[2pt]
\texttt{S15} & no rule for equal-size blocks \\\addlinespace[2pt]
\texttt{S16} & equations are (1) to (4) \\\addlinespace[2pt]
\texttt{F1} & numeric tags, no labels \\\addlinespace[2pt]
\texttt{F2} & Flock's fixed coordinate has no provenance \\\addlinespace[2pt]
\texttt{F3} & one root for push and pull \\\addlinespace[2pt]
\texttt{F4} & sumcheck target derived, not sent \\\addlinespace[2pt]
\texttt{F5} & bus forms share the last three powers \\\addlinespace[2pt]
\texttt{F6} & cubic round polynomial, three scalars \\\addlinespace[2pt]
\texttt{F7} & one coefficient never sent \\\addlinespace[2pt]
\texttt{F8} & ring-switched claims take the low powers \\\addlinespace[2pt]
\texttt{F9} & Python verifier omits four caps \\\addlinespace[2pt]
\texttt{F10} & rejection set of the Rust verifier \\\addlinespace[2pt]
\texttt{F11} & count tree holds the tables' counts only \\\addlinespace[2pt]
\texttt{F12} & fill blocks make heights exact \\\addlinespace[2pt]
\texttt{F13} & grinding binds the nonce on failure \\\addlinespace[2pt]
\texttt{F14} & seed hashes a constant naming the circuit \\\addlinespace[2pt]
\texttt{F15} & stack height bound checked separately \\\addlinespace[2pt]
\texttt{F16} & security level and no grinding for the bus \\\addlinespace[2pt]
\texttt{F17} & equal-size blocks in column-index order \\\addlinespace[2pt]
\texttt{F18} & one equation on the words, not one per limb \\
\end{longtable}}

\textbf{Findings about the libraries} (to \texttt{docs/\allowbreak{}dependencies.\allowbreak{}md}) and \textbf{about Lean} (to \texttt{tests/\allowbreak{}README.\allowbreak{}md}): the one-sentence descriptions of \cref{tab:docs-code-index} serve as names; the former codes \texttt{A10}-\texttt{A19}, \texttt{C5}, \texttt{C6}, \texttt{C10}, \texttt{C11}, \texttt{P1}, \texttt{P3}-\texttt{P6}, \texttt{E6}-\texttt{E15}, \texttt{E18} are kept in one column of each register. The two ``environment'' findings that are rules (\texttt{E16}, \texttt{E17}) are already the blueprint's convention \emph{Generic code} (\texttt{protocol-\allowbreak{}blueprint.\allowbreak{}md:\allowbreak{}329}) and are deleted.

\textbf{Findings of the three reviews}: the names are the third column of the review rows of \cref{tab:docs-code-index}. They are used inside the handoffs only.

\subsection{The index, drafted}\label{sec:docs-codes-index-draft}

One table, in the blueprint, section 13. It lists what keeps a number and what had a code.

{\footnotesize\setlength{\tabcolsep}{3pt}
\begin{longtable}{@{}>{\raggedright\arraybackslash}p{0.208\linewidth}>{\raggedright\arraybackslash}p{0.141\linewidth}>{\raggedright\arraybackslash}p{0.104\linewidth}>{\raggedright\arraybackslash}p{0.292\linewidth}>{\raggedright\arraybackslash}p{0.198\linewidth}@{}}
\caption{The index of names and numbers, drafted for the blueprint.}\label{tab:docs-index-draft}\\
\toprule
\textbf{Name} & \textbf{Number or former code} & \textbf{Kind} & \textbf{In one sentence} & \textbf{Specified or decided at} \\
\midrule
\endfirsthead
\toprule
\textbf{Name} & \textbf{Number or former code} & \textbf{Kind} & \textbf{In one sentence} & \textbf{Specified or decided at} \\
\midrule
\endhead
\bottomrule
\endfoot
base proof extraction and completeness & T4 & target theorem & an accepted proof has an execution, and an execution has an accepted proof & \texttt{architecture.\allowbreak{}md}, ``Target theorem ladder'' \\\addlinespace[2pt]
arithmetization and ISA equivalence & T1 & target theorem & constraint assignments and executions correspond, for well-formed programs & \texttt{architecture.\allowbreak{}md}; leanISA blueprint, Layer 10 \\\addlinespace[2pt]
witness-generator correctness & T2 & target theorem & the generated assignment satisfies the constraints; out of scope here & \texttt{architecture.\allowbreak{}md} \\\addlinespace[2pt]
recursive-verifier correctness & T5 & target theorem & out of scope here; \texttt{verify} is shaped for it & \texttt{architecture.\allowbreak{}md}; blueprint, ``Boundaries'' \\\addlinespace[2pt]
recursion extraction & T6 & target theorem & open research; out of scope here & \texttt{architecture.\allowbreak{}md} \\\addlinespace[2pt]
field instances & Layer 0; formerly hole \texttt{0} & unit of work & ArkLib as a dependency; sampling and oracle instances for the fields & blueprint, ``The units of work'' \\\addlinespace[2pt]
tables and stacking & Layer 1; formerly \texttt{L1} & unit of work & hypercube tables, stacking with selectors, the index and bytecode columns & same \\\addlinespace[2pt]
the spine & formerly \texttt{S} & unit of work & the instance, the relation, the seams, the phase interfaces, the composition & blueprint, ``The spine'' \\\addlinespace[2pt]
knowledge-soundness composition & formerly hole \texttt{C1} & unit of work & the port of ArkLib pull request 615, deleted when the pin passes it & blueprint, ``The spine'' \\\addlinespace[2pt]
Clean expressions as polynomials & Layer 2; formerly \texttt{I1} & unit of work & constraint and flush polynomials read off a Clean component & blueprint, ``The units of work'' \\\addlinespace[2pt]
the adaptor & Layer 3; formerly \texttt{I2} & unit of work & the leanISA instance and the two maps between the stack and Clean's witness & same \\\addlinespace[2pt]
sumcheck: definition and completeness & Layer 4; formerly \texttt{G1} & unit of work & \textbf{[TECHNICAL CHANGES TO BE INSERTED BY THE ORCHESTRATOR: one sentence, after the review of Layer 4]} & same \\\addlinespace[2pt]
sumcheck: knowledge soundness & Layer 4; formerly \texttt{G2} & unit of work & \textbf{[TECHNICAL CHANGES TO BE INSERTED BY THE ORCHESTRATOR]} & same \\\addlinespace[2pt]
batching by powers & Layer 4; formerly \texttt{G3} & unit of work & several claims become one by the powers of one challenge & same \\\addlinespace[2pt]
fingerprint and collision bound & Layer 5; formerly \texttt{G4} & unit of work & specification Lemma 5.2 and Theorem 5.1 & same \\\addlinespace[2pt]
grand-product GKR: definition and completeness & Layer 5; formerly \texttt{G5} & unit of work & \textbf{[TECHNICAL CHANGES TO BE INSERTED BY THE ORCHESTRATOR]} & same \\\addlinespace[2pt]
grand-product GKR: knowledge soundness & Layer 5; formerly \texttt{G6} & unit of work & \textbf{[TECHNICAL CHANGES TO BE INSERTED BY THE ORCHESTRATOR]} & same \\\addlinespace[2pt]
bus phase: definition and completeness & Layer 6; formerly \texttt{P1} & unit of work & \textbf{[TECHNICAL CHANGES TO BE INSERTED BY THE ORCHESTRATOR]} & same \\\addlinespace[2pt]
bus phase: knowledge soundness & Layer 6; formerly \texttt{P2} & unit of work & \textbf{[TECHNICAL CHANGES TO BE INSERTED BY THE ORCHESTRATOR]} & same \\\addlinespace[2pt]
table sumcheck phase: definition and completeness & Layer 7; formerly \texttt{P3} & unit of work & \textbf{[TECHNICAL CHANGES TO BE INSERTED BY THE ORCHESTRATOR]} & same \\\addlinespace[2pt]
table sumcheck phase: knowledge soundness & Layer 7; formerly \texttt{P4} & unit of work & \textbf{[TECHNICAL CHANGES TO BE INSERTED BY THE ORCHESTRATOR]} & same \\\addlinespace[2pt]
public-input phase & Layer 8; formerly \texttt{P5} & unit of work & specification section 8.2: one challenge, the sent values checked per limb, one claim per public line & same \\\addlinespace[2pt]
Flock phase & Layer 9; formerly \texttt{P6} & unit of work & the interface the Flock roadmap (issue 3) inhabits & same \\\addlinespace[2pt]
opening phase: definition and completeness & Layer 10; formerly \texttt{P7} & unit of work & \textbf{[TECHNICAL CHANGES TO BE INSERTED BY THE ORCHESTRATOR: pending the question of specification section 8.5]} & same \\\addlinespace[2pt]
opening phase: knowledge soundness & Layer 10; formerly \texttt{P8} & unit of work & \textbf{[TECHNICAL CHANGES TO BE INSERTED BY THE ORCHESTRATOR]} & same \\\addlinespace[2pt]
WHIR opening & Layer 11; formerly \texttt{K1} & unit of work & \textbf{[TECHNICAL CHANGES TO BE INSERTED BY THE ORCHESTRATOR]} & same \\\addlinespace[2pt]
Merkle trees, byte hasher, WHIR parameters & Layer 11; formerly \texttt{K2} & unit of work & \textbf{[TECHNICAL CHANGES TO BE INSERTED BY THE ORCHESTRATOR: see \cref{sec:docs-merkle} on VCVio's Merkle trees]} & same \\\addlinespace[2pt]
transcript, proof object, compiled verifier & Layer 12; formerly \texttt{K3} & unit of work & \textbf{[TECHNICAL CHANGES TO BE INSERTED BY THE ORCHESTRATOR]} & same \\\addlinespace[2pt]
the base theorems & Layer 13; formerly \texttt{K4} & unit of work & \textbf{[TECHNICAL CHANGES TO BE INSERTED BY THE ORCHESTRATOR: see \cref{sec:docs-leanisa} on the program hypothesis]} & same \\\addlinespace[2pt]
heights are powers of two & decision 1 & decision, taken & leanISA's \texttt{Caps} requires it, so the two roadmaps state one relation & leanISA blueprint, Layer 8; issue 13 \\\addlinespace[2pt]
statement and parameters & decision 2 & decision, taken & the public input is the statement; the instance indexes the protocol family & convention \emph{Statements and parameters} \\\addlinespace[2pt]
where the zerocheck is charged & decision 3 & decision, taken & in the bus phase, where the point is drawn; no separate phase & convention \emph{Seams} \\\addlinespace[2pt]
where generic code lives & decision 4 & decision, open & the \texttt{To<Library>} folders, or an ArkLib branch & blueprint, ``Decisions'' \\\addlinespace[2pt]
an end-to-end run of the honest prover & decision 5 & decision, open & depends on the cost of the encoder & blueprint, ``Decisions'' \\\addlinespace[2pt]
the witness is the stack & decision 6 & decision, taken & the protocol's witness is the committed column, in \texttt{Type 0} & ``The relation ladder'' \\\addlinespace[2pt]
phases over an abstract instance & decision 7 & decision, taken & every phase is tested on the toy instance & convention \emph{The wall} \\\addlinespace[2pt]
five clauses, caps outside & decision 8 & decision, taken & the caps are a hypothesis of the adaptor's soundness theorem & ``What the spine fixes'' \\\addlinespace[2pt]
balance is a permutation & decision 9 & decision, taken & the pushed and pulled tuple lists are permutations of each other & ``What the spine fixes'' \\\addlinespace[2pt]
worst-case round-by-round knowledge soundness & decision 10 & decision, taken & per component, composed by the ported theorem & convention \emph{Holes} \\\addlinespace[2pt]
extractors are computable definitions & decision 11 & decision, taken & the protocol's extractor is \texttt{piopExtractor} & convention \emph{Extractors}; acceptance test 24 \\\addlinespace[2pt]
the strong Flock predicate & decision 12 & decision, taken & the predicate is Flock's R1CS on the committed region & Layers 3 and 9 \\\addlinespace[2pt]
public lines, not cells & decision 13 & decision, taken & the statement fixes cells 0 and 1 of listed columns & convention \emph{Public input} \\\addlinespace[2pt]
the degree bound belongs to the instance & decision 14 & decision, taken & the bus seam bounds every term by it & convention \emph{Seams}; acceptance test 28 \\\addlinespace[2pt]
the public-input transcript & decision 15 & decision, taken & the prover sends the values; a line says whether its value is sent & convention \emph{Public input}; Layer 8 \\\addlinespace[2pt]
seams carry claims, never challenges & none today & decision, taken & what a phase emits is fixed by its knowledge soundness & ``What the spine fixes'' \\\addlinespace[2pt]
schedules travel with the definitions & none today & decision, taken & the spine fixes the commit phase's schedule only & ``What the spine fixes'', item 4 \\\addlinespace[2pt]
fingerprint degree & acceptance test 1 & acceptance test & the error carries a factor 4 & blueprint, ``Acceptance tests'' \\\addlinespace[2pt]
padding leaves are 1 & acceptance test 2 & acceptance test & a zero pad zeroes every product & same \\\addlinespace[2pt]
the nonzero count root is load-bearing & acceptance test 3 & acceptance test & a read with count zero balances & same \\\addlinespace[2pt]
one root, not two & acceptance test 4 & acceptance test & the push and pull roots are one scalar & same \\\addlinespace[2pt]
domain separators & acceptance test 5 & acceptance test & coordinate 0 separates the three buses & same \\\addlinespace[2pt]
zerocheck point recycling & acceptance test 6 & acceptance test & the point is drawn after the commitment & same \\\addlinespace[2pt]
back-loaded padding & acceptance test 7 & acceptance test & the lift keeps the table's sum & same \\\addlinespace[2pt]
degree three, three scalars & acceptance test 8 & acceptance test & the wire carries three coefficients & same \\\addlinespace[2pt]
shared bus powers & acceptance test 9 & acceptance test & the three bus forms share the last three powers & same \\\addlinespace[2pt]
top limb of the public input & acceptance test 10 & acceptance test & the third claim is pooled though nothing is sent for it & same \\\addlinespace[2pt]
joint list binding & acceptance test 11 & acceptance test & one member of the list satisfies all claims & same \\\addlinespace[2pt]
absorb before squeeze & acceptance test 12 & acceptance test & sizes and root are observed before the first challenge & same \\\addlinespace[2pt]
index column bit order & acceptance test 13 & acceptance test & bit \texttt{k} is coordinate \texttt{k} & same \\\addlinespace[2pt]
bytecode slot bits & acceptance test 14 & acceptance test & the opcode is slot 3, low bit first; witness \texttt{bytecode\allowbreak{}Column\_\allowbreak{}answer\_\allowbreak{}bool\allowbreak{}Vec} & same \\\addlinespace[2pt]
selector alignment & acceptance test 15 & acceptance test & blocks sit at multiples of their size, largest first & same \\\addlinespace[2pt]
variable order & acceptance test 16 & acceptance test & the highest variable is bound first & same \\\addlinespace[2pt]
radix and parity & acceptance test 17 & acceptance test & one radix-2 layer when the height is odd & same \\\addlinespace[2pt]
counts in the count tree & acceptance test 18 & acceptance test & the tables' count columns only & same \\\addlinespace[2pt]
the extractor reads the stack & acceptance test 19 & acceptance test & the witness is read off the committed column & same \\\addlinespace[2pt]
no exceptional challenge & acceptance test 20 & acceptance test & perfect completeness has no side condition & same \\\addlinespace[2pt]
sizes are parameters & acceptance test 21 & acceptance test & the protocol is a family indexed by the sizes & same \\\addlinespace[2pt]
tags, not labels & acceptance test 22 & acceptance test & four numeric tags & same \\\addlinespace[2pt]
numeric error & acceptance test 23 & acceptance test & the bound at the caps & same \\\addlinespace[2pt]
the extractor computes & acceptance test 24 & acceptance test & an extractor chosen classically proves soundness, not knowledge & same \\\addlinespace[2pt]
the import rule holds & acceptance test 25 & acceptance test & no phase imports the arithmetization & same \\\addlinespace[2pt]
seams are the contract & acceptance test 26 & acceptance test & a phase's output relation is the next one's input relation by definition & same \\\addlinespace[2pt]
the toy instance is honest & acceptance test 27 & acceptance test & each clause fails alone & same \\\addlinespace[2pt]
the bus seam bounds the degree & acceptance test 28 & acceptance test & a cubic term is rejected & same \\\addlinespace[2pt]
admitted upstream: sumcheck round knowledge soundness & formerly ledger \texttt{A1} & ledger entry & the single-round bound is admitted in ArkLib at the pin & blueprint, ``Dependencies'', the ledger \\\addlinespace[2pt]
admitted upstream: composition of knowledge soundness & formerly ledger \texttt{A2} & ledger entry & proved here by the port; the port is deleted at the pin bump & same \\\addlinespace[2pt]
(the other eleven ledger entries, as in \cref{sec:docs-codes-mapping}) &  & ledger entry &  & same \\
\end{longtable}}

\textbf{[TECHNICAL CHANGES TO BE INSERTED BY THE ORCHESTRATOR: rows for the findings against the sources that the blueprint cites, once the other agents' findings are merged into the register.]}
